\chapter{复习路线与自检题}

\section{遇到问题时先选哪种方法}
\begin{center}
\small
\begin{tabular}{p{.24\textwidth}p{.29\textwidth}p{.36\textwidth}}
\hline
问题的物理结构 & 优先工具 & 先检查的条件\\
\hline
少量质点、受力明确 & Newton方程与受力图 & 惯性系、固定质量、接触状态\\
只问速度与位置的关系 & 功能关系、能量 & 外部驱动、耗散、移动约束是否做功\\
孤立两体、中心相互作用 & 质心分离、角动量、有效势 & 势只依赖距离；轨道能量是哪一个系的量\\
完整理想约束较多 & 广义坐标、Lagrange方程 & 独立自由度；坐标消元是否合法\\
稳定平衡附近耦合运动 & 质量矩阵与刚度矩阵 & 小振幅、非退化性、有无速度耦合\\
刚体自由转动 & 随体主轴、Euler方程 & 惯量参考点、力矩所在坐标系\\
周期运动、参数缓变 & 作用量--角变量 & 远离分离曲线；变化慢于周期\\
长期数值轨道 & 结构保持积分与收敛检查 & 适用的Hamiltonian形式及步长稳定性\\
\hline
\end{tabular}
\end{center}

\section{把常见误用改成条件明确的陈述}
能量、动量和角动量不是看到题目就同时可用的三个常数。
机械能需核对功率，动量需核对系统边界和外力，角动量需核对取矩点和外力矩。
旋转系中的有效势、Lagrange能量函数以及惯性系机械能常常不同。

小振动频率来自$\det(K-\omega^2M)=0$，不是单独对$K$取特征值；
阻尼振子的自由衰减频率、位移共振频率与功率峰频率也不相同。
半正定势能Hessian、零频率模和临界参数都应继续分析高阶项或对称性。

正则动量来自$\p L/\p\dot q$，不应一律写为$mv$。
Hamiltonian自治才保证它自身守恒，辛积分也不等于精确能量守恒。
守恒量足以降低方程维数，但一个能量常数通常不能保证多自由度系统可积。

\section{短题与核对结果}
下列短题以复算关键步骤为目的。先遮住核对结果完成推导，再检查条件和极限。

\subsection*{1. 圆轨道的两个频率}
质量为$\mu$的粒子在势能$V(r)=ar^4$中运动，$a>0$。给定非零角动量$\ell$，求圆轨道半径及两种小运动频率。

\paragraph{核对}
由$4ar_0^3=\ell^2/(\mu r_0^3)$得$r_0=(\ell^2/4\mu a)^{1/6}$。
$\Omega^2=4ar_0^2/\mu$，$\omega_r^2=24ar_0^2/\mu=6\Omega^2$。
相邻近心点之间的方位角约为$2\pi/\sqrt6$，不能预先假定轨道是闭合椭圆。

\subsection*{2. 共振时差一个什么因子}
阻尼振子满足$\ddot x+2\gamma\dot x+\omega_0^2x=(F_0/m)\cos\omega_0t$，$\gamma>0$。
求稳态振幅、相位滞后和平均耗散功率。

\paragraph{核对}
$A=F_0/(2m\gamma\omega_0)$，$\delta=\pi/2$，
$\ave P=m\gamma\omega_0^2A^2=F_0^2/(4m\gamma)$。
检查阻尼项写的是$2\gamma\dot x$还是$\gamma\dot x$，可避免整套结果相差2倍。
$\gamma\to0$不能在固定“已达稳态”的假设下直接取极限。

\subsection*{3. 竖直加速的电梯中的摆}
电梯支点高度为$a(t)$，竖直向上为正，推导长度$l$的单摆方程。

\paragraph{核对}
$x=l\sin\theta$，$y=a-l\cos\theta$给出
$L=ml^2\dot\theta^2/2+ml\dot a\dot\theta\sin\theta+mgl\cos\theta$，
略去仅含时间的项后可得
$\ddot\theta+(g+\ddot a)\sin\theta/l=0$。
向上加速增大有效重力，自由落体$\ddot a=-g$时重力回复项消失。
若电梯加速度随时间变化，应先判断能否使用瞬时频率近似。

\subsection*{4. 自由链的零模}
两个质量分别为$m_1,m_2$的质点沿同一直线运动，仅用劲度$k$的弹簧连接，没有固定墙。
求简正频率并解释零模。

\paragraph{核对}
以$x_1,x_2$表示相对选定平衡构型的位移，$M=\diag(m_1,m_2)$，$K=k\left(\begin{smallmatrix}1&-1\\-1&1\end{smallmatrix}\right)$。
$\omega_0=0$对应整体平移，非零频率为
$\omega_1^2=k(1/m_1+1/m_2)=k/\mu$。
非零模满足$m_1x_1+m_2x_2=0$。零模的速度初值产生质心匀速漂移，
并非总绕某个固定平衡位置振动。

\subsection*{5. 不滑动能否实现}
实心均匀圆柱从静止沿倾角为$\alpha$的斜面滚下，$I=\beta mR^2$、$\beta=1/2$。求加速度与最小静摩擦系数。

\paragraph{核对}
$a=2g\sin\alpha/3$，$f=mg\sin\alpha/3$，
故$\mu_s\geq\tan\alpha/3$。
若不满足此不等式，必须改写为含滑动摩擦的平移、转动两个方程，
不能继续强加$a=R\dot\omega$。

\subsection*{6. 有磁场时哪个动量守恒}
质量为$m$、电荷为$e$的粒子处于均匀磁场$\bm B=B\nvec z$中，取$\bm A=(0,Bx,0)$，无电场。
写出Hamiltonian并指出一个循环坐标的动量。

\paragraph{核对}
$H=[p_x^2+(p_y-eBx)^2+p_z^2]/(2m)$，$y$不出现，故$p_y$守恒。
机械动量$mv_y=p_y-eBx$随$x$变化，一般不守恒；$p_z=mv_z$也守恒。
规范变换改变正则动量的表达，但不会改变实际圆周运动和沿磁场的匀速运动。

\subsection*{7. 骤变与缓变的区别}
一个质量固定的谐振子频率从$\omega_i$变为$\omega_f$。
在平滑绝热改变与瞬间改变两种情形下，最终能量怎样变化？

\paragraph{核对}
在第9章的慢时间假设下，若$\omega(t)>0$且远离零，绝热近似给出$E_f\simeq E_i\omega_f/\omega_i$。
仅有$|\dot\omega|/\omega^2\ll1$不足以排除参量共振。
瞬间改变刚度且不施加冲量时，$q,p$连续，故
$E_f=E_i+\tfrac12m(\omega_f^2-\omega_i^2)q^2$，依赖改变瞬间的相位。
对初始相位均匀平均后，$\ave{E_f}=\tfrac12E_i[1+(\omega_f/\omega_i)^2]$。
这两种过程不应使用同一守恒假设。

\subsection*{8. 一个缩放是否正则}
令$Q=aq$、$P=bp$，其中$a,b$为非零常数。判断何时正则，并给出生成函数。

\paragraph{核对}
$\{Q,P\}=ab$，故正则的充要条件为$ab=1$。
取$F_2(q,P)=aqP$即可得到$Q=aq$、$p=aP$。
仅把坐标与动量同时放大相同倍数一般不正则；一自由度相面积须保持。

\subsection*{9. Verlet的稳定界限}
对$V=m\omega_0^2q^2/2$，将速度Verlet写成$(q_{n+1},p_{n+1})$关于$(q_n,p_n)$的线性映射。

\paragraph{核对}
取$h>0$、$\omega_0>0$，令$u=h\omega_0$，矩阵为
\begin{equation}
 S_h=\begin{pmatrix}
 1-u^2/2 & h/m\\
 -hm\omega_0^2(1-u^2/4)&1-u^2/2
 \end{pmatrix}.
\end{equation}
$\det S_h=1$对任意步长成立，但特征值只有在$0<u<2$时位于单位圆上且不退化。
$u=2$时$S_h=\left(\begin{smallmatrix}-1&h/m\\0&-1\end{smallmatrix}\right)$，是非平凡Jordan块，因而
\[
 q_n=(-1)^n\left(q_0-\frac{nh}{m}p_0\right),\qquad p_n=(-1)^np_0.
\]
除$p_0=0$的特殊初态外，位置振幅线性增长；$u>2$时一般初态则有指数增长。
这说明保相面积并不能代替稳定性检查。
